\documentclass[11pt,a4paper]{article}
\usepackage{turkos}

\begin{document}
\renewcommand{\progresslabel}{Stage V progress toward Turk-OS 2.0 (this phase in red)}
\phaseheader{17}{The RISC-V Port}{Boot the same kernel and userland on a second architecture: QEMU's RISC-V virt board, with OpenSBI as the firmware and Turk-OS in supervisor mode}{20}{60}

\ataglance{5--6 weeks (Task 17.13 on your own core is open-ended)}{4}{Phase 16 complete: the portable core, both cross
compilers, the FDT parser and its host tests}{Turk-OS on \code{qemu-system-riscv32 -machine virt}, booting to
\code{tsh}: traps with a crash report that names every cause, the SBI timer at 100\,Hz, the PLIC with UART input, an
Sv32 higher-half kernel at \code{0xC0000000}, user mode through \code{ecall}, a virtio-blk disk carrying TurkFS, and
\code{make ARCH=riscv32 test}. Optionally, a checked plan for running it on Timur-RV32IMC.}

\section{Why this phase}

Phase 16 claimed that everything outside \code{kernel/arch/i386} is generic. This phase tests the claim on a second
machine: you write the RISC-V half of \code{include/arch.h}, and the scheduler, the VFS, TurkFS, \code{fork} and the
whole userland are recompiled, not rewritten. Every place where that fails is a leak the audit of Task~16.2 missed; fix
it in the generic code, note it in \code{docs/porting.md}, and keep \code{make ARCH=i386 test} passing.

RISC-V comes first because it changes little at once. RV32 keeps the 32-bit word, and Sv32 is two-level 4\,KiB paging
like x86's, so the kernel stays at \code{0xC0000000} and Phase 10's user layout stays as it is. What changes is the boot
path (firmware instead of GRUB), traps (CSRs instead of an IDT and a TSS), devices (a devicetree instead of fixed PC
addresses) and the assembler (GNU \code{as} instead of NASM). And the road leads on to your own core (Task~17.13).

\section{Concepts you need}

\subsection{RISC-V in one page}

RISC-V is a load/store architecture: only loads and stores touch memory (\code{lw a0, 4(a1)} is x86's
\code{mov eax, [ebx+4]}), and there is no flags register, since branches compare two registers (\code{blt a0, a1, label}).
The registers \code{x0}--\code{x31} go by ABI names: \code{zero} (always 0), \code{ra} (return address), \code{sp},
\code{gp} and \code{tp} (unused by the kernel), \code{t0}--\code{t6} (temporaries), \code{s0}--\code{s11}
(callee-saved) and \code{a0}--\code{a7} (arguments; \code{a0}, \code{a1} also results). \code{call} leaves the return
address in \code{ra} instead of pushing it, and \code{sp} must be 16-byte aligned at every call.

The ISA string \code{rv32imac_zicsr_zifencei} means: 32-bit base integer (I), multiply and divide (M), atomics (A),
compressed 16-bit encodings (C), the CSR instructions (Zicsr) and \code{fence.i} (Zifencei). With C, assembler and
compiler choose 2-byte forms on their own, so an instruction is 2 or 4 bytes long; a 4-byte one has binary \code{11} in
its lowest two bits. \code{-mabi=ilp32} means 32-bit \code{int}, \code{long} and pointers and no floating-point
registers.

\subsection{Three privilege modes and the SBI}

A RISC-V hart (hardware thread) runs in machine (M), supervisor (S) or user (U) mode. Turk-OS runs in S-mode,
programs in U-mode. Unlike the BIOS and GRUB, the firmware, \textbf{OpenSBI}, stays resident in M-mode and offers
services through the \textbf{Supervisor Binary Interface} (SBI), called with \code{ecall} one level down
(Figure~\ref{fig:modes}).

\begin{figure}[htbp]
\centering
\begin{tikzpicture}[lay/.style={draw=tkNavy!60, rounded corners=2pt, minimum width=5.4cm, text width=5.2cm,
    minimum height=0.72cm, font=\sffamily\scriptsize, align=center, line width=0.6pt},
  arr/.style={-{Stealth[length=2mm]}, line width=0.7pt, draw=tkViolet}]
  \node[font=\sffamily\small\bfseries, text=tkNavy] at (0,0.7) {Turk-OS 1.0 on x86};
  \node[lay, fill=tkGreen!10] (xu) at (0,0) {ring 3: \code{tsh}, coreutils, Turk-libc};
  \node[lay, fill=tkRed!8, below=0.5cm of xu] (xk) {ring 0: the Turk-OS kernel};
  \node[lay, fill=tkGray!12, below=0.5cm of xk] (xf) {BIOS and GRUB: load the kernel, then gone};
  \node[lay, fill=tkGray!20, below=0.5cm of xf] (xh) {PC hardware};
  \node[font=\sffamily\small\bfseries, text=tkViolet] at (7.6,0.7) {Turk-OS on RISC-V \code{virt}};
  \node[lay, fill=tkGreen!10] (ru) at (7.6,0) {U-mode: the same programs, rebuilt};
  \node[lay, fill=tkRed!8, below=0.5cm of ru] (rk) {S-mode: the Turk-OS kernel};
  \node[lay, fill=tkViolet!12, draw=tkViolet, below=0.5cm of rk] (rf) {M-mode: OpenSBI, stays resident};
  \node[lay, fill=tkGray!20, below=0.5cm of rf] (rh) {hart, ACLINT timer, PLIC, 16550, virtio-mmio};
  \draw[arr] ([xshift=1.2cm]ru.south) -- node[right, font=\sffamily\scriptsize, text=tkViolet] {\code{ecall}: system call} ([xshift=1.2cm]rk.north);
  \draw[arr] ([xshift=1.2cm]rk.south) -- node[right, font=\sffamily\scriptsize, text=tkViolet] {\code{ecall}: SBI call} ([xshift=1.2cm]rf.north);
  \draw[tkarrow] ([xshift=1.2cm]xu.south) -- node[right, font=\sffamily\scriptsize] {\code{int 0x80}} ([xshift=1.2cm]xk.north);
  \draw[tkarrow, dashed] ([xshift=-1.2cm]xf.north) -- node[left, font=\sffamily\scriptsize] {loads} ([xshift=-1.2cm]xk.south);
  \draw[tkarrow, dashed] ([xshift=-1.2cm]rf.north) -- node[left, font=\sffamily\scriptsize] {\code{sret} into} ([xshift=-1.2cm]rk.south);
\end{tikzpicture}
\caption{Where the code runs. OpenSBI is to S-mode what the BIOS and GRUB were to ring 0, except that it never leaves.}
\label{fig:modes}
\end{figure}

An SBI call puts the extension ID (EID) in \code{a7}, the function ID (FID) in \code{a6} and arguments in
\code{a0}--\code{a5}, and returns an error code in \code{a0} and a value in \code{a1}; other registers are preserved.
Turk-OS uses Base (EID \code{0x10}, to probe), Timer (\code{0x54494D45}, ``TIME''), System Reset (\code{0x53525354},
``SRST'') and Debug Console (\code{0x4442434E}, ``DBCN''), with the legacy putchar (EID \code{0x01}) as a fallback.
OpenSBI v1.7, QEMU 10.2's \code{-bios default}, lists all four in its banner.

\subsection{Control and status registers}

What x86 did with special registers and tables, RISC-V does with \textbf{control and status registers} (CSRs), read
and written by \code{csrr}, \code{csrw}, \code{csrs} (set bits), \code{csrc} (clear bits) and \code{csrrw} (swap). The
supervisor CSRs of this phase, with their RV32 bit positions from the privileged specification:
\code{sstatus} (bit 1 SIE, interrupts on, like EFLAGS.IF; bit 5 SPIE, SIE before the trap; bit 8 SPP, the mode before
the trap, 0 = U; bits 14:13 FS, the FPU state; bit 18 SUM, S-mode may touch U pages); \code{sie} and \code{sip} (bit 1
software, 5 timer, 9 external interrupts); \code{stvec} (the trap handler, like the IDT); \code{sscratch} (free for the
kernel, like the TSS's \code{esp0}); \code{sepc}, \code{scause} and \code{stval} (the trapping pc, the cause with bit
31 set for interrupts, and the faulting address like CR2); \code{satp} (bit 31 MODE, 1 = Sv32; bits 30:22 ASID; bits
21:0 the root table's PPN, like CR3); and the read-only 64-bit \code{time} counter (\code{rdtime}, \code{rdtimeh}).

\subsection{What happens on a trap}

Exceptions, interrupts and system calls all jump to \code{stvec}. The hardware writes \code{sepc}, \code{scause} and
\code{stval}, moves the old mode to SPP and SIE to SPIE, clears SIE and jumps (Figure~\ref{fig:trap}). It switches no
stack and saves no register; \code{entry.S} does that, and \code{sret} undoes the CSR part. \code{sepc} points
\emph{at} the trapping instruction, which is right for a page fault; after \code{ecall} and \code{ebreak} the kernel
must advance it, or the program traps forever.

\begin{figure}[htbp]
\centering
\begin{tikzpicture}[node distance=0.35cm and 0.45cm, every node/.style={font=\sffamily\scriptsize}]
  \node[tkok, text width=2.75cm] (u) {\textbf{U-mode}: \code{write}\\a7 = 4, a0--a2 = arguments\\\code{ecall}};
  \node[tkbox, text width=3.15cm, right=of u] (hw) {\textbf{hardware}: sepc = pc; scause = 8; stval = 0; SPP = 0; SPIE = SIE; SIE = 0; pc = stvec};
  \node[tkbox, text width=3.0cm, right=of hw] (en) {\textbf{entry.S}: swap sp with \code{sscratch}; save 31 registers and 4 CSRs};
  \node[tkhi, text width=3.0cm, below=0.5cm of en] (c) {\textbf{trap.c}: cause 8, so sepc += 4; enable SIE; \code{syscall_dispatch}};
  \node[tkbox, text width=3.15cm, left=of c] (ret) {\textbf{trap\_return}: restore; \code{sscratch} = top of the kernel stack};
  \node[tkok, text width=2.75cm, left=of ret] (s) {\code{sret}: pc = sepc, mode = SPP, SIE = SPIE; a0 = result};
  \draw[tkarrow] (u) -- (hw); \draw[tkarrow] (hw) -- (en); \draw[tkarrow] (en) -- (c);
  \draw[tkarrow] (c) -- (ret); \draw[tkarrow] (ret) -- (s);
\end{tikzpicture}
\caption{A system call on RISC-V. Compare Figure 1 of Phase 10: no gate, no TSS, no stack switch in hardware.}
\label{fig:trap}
\end{figure}

\code{scause} replaces the x86 vector (Table~\ref{tab:causes}). Note that division by zero does not trap at all,
and that a page fault comes without an error code: whether the page was missing or protected, the kernel finds out by
looking it up.

\begin{table}[htbp]
\centering\small\renewcommand{\arraystretch}{1.15}
\begin{tabularx}{\linewidth}{@{}>{\raggedright\arraybackslash}p{0.25\linewidth} >{\raggedright\arraybackslash}p{0.36\linewidth} X@{}}
\toprule
\tkth{x86 (Phases 3, 4, 10)} & \tkth{RISC-V \code{scause}} & \tkth{Notes} \\
\midrule
\#DE divide error (0) & none & the quotient is all ones \\
\#BP \code{int3} (3) & 3 breakpoint (\code{ebreak}) & \code{sepc} points at the \code{ebreak} \\
\#UD (6), \#GP (13) & 2 illegal instruction; 5 load, 7 store/AMO access fault & also a privileged CSR from U-mode \\
\#PF (14), CR2 & 12 instruction, 13 load, 15 store/AMO page fault & address in \code{stval} \\
\#AC (alignment) & 0, 4, 6 address misaligned & OpenSBI emulates loads and stores \\
\code{int 0x80} & 8 ecall from U-mode & 9: from S-mode, OpenSBI's \\
IRQ 0, the PIT & interrupt 5, supervisor timer & armed through the SBI \\
IRQ 1, 4, 14 & interrupt 9, supervisor external & ask the PLIC which device \\
\bottomrule
\end{tabularx}
\caption{x86 vectors and RISC-V causes. An interrupt has bit 31 set: a timer interrupt reads \code{0x80000005}.}
\label{tab:causes}
\end{table}

\subsection{Interrupts: the timer through the SBI, devices through the PLIC}

S-mode sees three interrupt sources, each enabled by a bit in \code{sie} and all of them by \code{sstatus}.SIE. The
\textbf{timer} comparator belongs to M-mode on this board, so you arm it with an SBI call; it fires once, and the
handler sets the next deadline. The \code{time} counter runs at the devicetree's \code{timebase-frequency}, 10\,MHz
here. \textbf{External} interrupts all arrive as cause 9: the \textbf{PLIC} (platform-level interrupt controller)
collects the device lines, and the handler reads its \emph{claim} register to learn the source and writes the number
back to \emph{complete} it, like the 8259's acknowledge and EOI.

\subsection{Sv32: two-level paging, the RISC-V way}

Sv32 walks the same tree as x86 (Figure~\ref{fig:sv32}): \code{satp} names a root table of 1024 four-byte
entries, VPN[1] selects one, VPN[0] an entry of the second-level table. A root entry can also map a 4\,MiB
\textbf{megapage}. Root entries 768--1023 hold the kernel half, shared by every address space as in Phase 10.

\begin{figure}[htbp]
\centering
\begin{tikzpicture}[x=0.4cm, y=0.9cm, every node/.style={font=\sffamily\scriptsize}]
  \draw[draw=tkNavy!70, fill=tkBlue!12] (0,4) rectangle (10,4.7);
  \draw[draw=tkNavy!70, fill=tkGreen!12] (10,4) rectangle (20,4.7);
  \draw[draw=tkNavy!70, fill=tkAmber!14] (20,4) rectangle (32,4.7);
  \node at (5,4.35) {VPN[1] (bits 31--22)};
  \node at (15,4.35) {VPN[0] (bits 21--12)};
  \node at (26,4.35) {offset (bits 11--0)};
  \node[font=\sffamily\scriptsize\bfseries, anchor=east] at (-0.3,4.35) {virtual address};
  \node[draw=tkViolet, fill=tkViolet!12, rounded corners=2pt, align=center, minimum width=1.6cm] (satp) at (-1.5,1.3) {\code{satp}\\PPN};
  \node[draw=tkNavy!70, fill=tkBlue!6, minimum width=2.2cm, minimum height=2.0cm, align=center] (pd) at (5,1.3) {root table\\(1024 PTEs)\\pointer, or a\\4\,MiB leaf};
  \node[draw=tkNavy!70, fill=tkGreen!6, minimum width=2.2cm, minimum height=2.0cm, align=center] (pt) at (15,1.3) {second-level\\table\\(1024 PTEs)};
  \node[draw=tkNavy!70, fill=tkAmber!8, minimum width=2.2cm, minimum height=2.0cm, align=center] (fr) at (26,1.3) {4\,KiB\\physical\\page};
  \draw[tkarrow] (satp) -- (pd.west);
  \draw[tkarrow] (5,4) -- node[right] {selects PTE} (pd.north);
  \draw[tkarrow] (15,4) -- node[right] {selects PTE} (pt.north);
  \draw[tkarrow] (26,4) -- node[right] {byte in page} (fr.north);
  \draw[tkarrow] (pd.east) -- node[above, align=center, font=\sffamily\tiny] {R=W=X=0:\\pointer} (pt.west);
  \draw[tkarrow] (pt.east) -- node[above, align=center, font=\sffamily\tiny] {leaf:\\page} (fr.west);
\end{tikzpicture}
\caption{Sv32 translation with 4\,KiB pages. Apart from the names, this is Figure 1 of Phase 6.}
\label{fig:sv32}
\end{figure}

The entries differ more (Table~\ref{tab:pte}). V (valid) replaces P; an entry with R, W and X all zero points to
the next table, any other is a \textbf{leaf}; every page has an execute bit. And a core may raise a page fault instead
of setting the \textbf{A} and \textbf{D} bits. QEMU sets them, so the bug hides there, but a simple core, possibly your
own, will not. Turk-OS never reads them, so it sets both in every leaf, as the specification recommends.

\begin{table}[htbp]
\centering\small\renewcommand{\arraystretch}{1.15}
\begin{tabularx}{\linewidth}{@{}>{\raggedright\arraybackslash}p{0.21\linewidth} >{\raggedright\arraybackslash}p{0.33\linewidth} X@{}}
\toprule
\tkth{Question} & \tkth{x86 PDE/PTE (Phase 6)} & \tkth{Sv32 PTE} \\
\midrule
In use? & P, bit 0 & V, bit 0 \\
Table or page? & PS bit in a PDE (4\,MiB page) & leaf if any of R, W, X is set \\
Permissions & R/W bit 1; no execute bit & R, W, X bits 1--3; W needs R \\
User access & U/S bit 2, in PDE \emph{and} PTE & U bit 4, in the leaf; S-mode needs \code{sstatus}.SUM to touch it \\
Accessed, dirty & A bit 5, D bit 6, set by the CPU & A bit 6, D bit 7; may fault instead \\
Free for the OS & bits 9--11 & RSW, bits 8--9 (Turk-OS: bit 8 = COW) \\
Frame number & bits 31:12, 32-bit physical & bits 31:10, 34-bit physical \\
Flush one page & \code{invlpg} & \code{sfence.vma va, zero} \\
\bottomrule
\end{tabularx}
\caption{The same questions, answered by two page-table formats. G (global) is bit 8 on x86, bit 5 on Sv32.}
\label{tab:pte}
\end{table}

\subsection{The virt board}

QEMU's \code{virt} board is a small, clean set of devices whose addresses and interrupts OpenSBI passes to the
kernel in a \textbf{flattened devicetree} (FDT), read by your Task~16.10 parser. Table~\ref{tab:virt} shows QEMU 10.2
with the default 128\,MiB; always take the real values from the FDT.

\begin{table}[htbp]
\centering\small\renewcommand{\arraystretch}{1.15}
\begin{tabularx}{\linewidth}{@{}l l >{\raggedright\arraybackslash}X l@{}}
\toprule
\tkth{Address} & \tkth{Size} & \tkth{Device and \code{compatible} string} & \tkth{PLIC IRQ} \\
\midrule
\code{0x00100000} & 4\,KiB & test device, \code{sifive,test0}: ends QEMU with a status & -- \\
\code{0x02000000} & 64\,KiB & CLINT (ACLINT timer and software interrupts), \code{riscv,clint0}; M-mode only & -- \\
\code{0x0C000000} & 6\,MiB & PLIC, \code{riscv,plic0}, 95 sources & -- \\
\code{0x10000000} & 256\,B & UART, \code{ns16550a}, 3.6864\,MHz clock & 10 \\
\code{0x10001000}--\code{0x10008000} & 4\,KiB each & eight virtio-mmio slots, \code{virtio,mmio} & 1--8 \\
\code{0x80000000} & 128\,MiB & RAM; OpenSBI keeps the first 320\,KiB & -- \\
\bottomrule
\end{tabularx}
\caption{The RISC-V \code{virt} board of QEMU 10.2. \code{/cpus} has \code{timebase-frequency = 10000000} and
\code{mmu-type = "riscv,sv32"}; \code{/chosen} has \code{stdout-path = "/soc/serial@10000000"}.}
\label{tab:virt}
\end{table}


\begin{decisionbox}{Run at the physical address first, move to 0xC0000000 in Task 17.7}
You could link at \code{0xC0400000} from day one and enable Sv32 in \code{boot.S}. Don't, yet: with paging off,
each bug in Tasks 17.2--17.6 lives in one new thing (the entry code, a CSR, a device address). Link at the physical
address OpenSBI jumps to, keep \code{satp} at 0, and make \code{P2V}/\code{V2P} the identity by setting
\code{KERNEL_BASE} to \code{PHYS_BASE} (\code{0x80000000}) in \code{<arch/memlayout.h>}. Task~17.7 then changes one
number and the linker script and adds a boot page table, and everything above keeps working because it only used the
macros: Phases 1--6 of x86, compressed into one phase.
\end{decisionbox}

\section{Reading plan}

\begin{readingplan}
\must & \RVSPEC & Privileged ISA (Vol.~II), chapter \emph{Supervisor-Level ISA}: \emph{Supervisor CSRs}, \emph{Supervisor Instructions}, \emph{Sv32} & -- & Every bit, cause and PTE flag of Tasks 17.4--17.9. \\
\must & \RVSPEC & SBI specification: ch.~3 \emph{Binary Encoding}, ch.~4 \emph{Base}, ch.~6 \emph{Timer}, ch.~10 \emph{System Reset}, ch.~12 \emph{Debug Console Extension} & -- & The firmware calls of Tasks 17.2, 17.5, 17.12. \\
\should & \RVSPEC & Unprivileged ISA (Vol.~I): \emph{RV32I Base Integer Instruction Set}, \emph{Zicsr}; skim \emph{C} and \emph{A} & -- & What you write by hand and read in GDB. \\
\must & \XVSIX & Ch.~4 \emph{Traps and system calls} & -- & The same trap path on RV64, line by line. \\
\should & \XVSIX & Ch.~3 \emph{Page tables}; Ch.~5 \emph{Page faults}; Ch.~6 \emph{Interrupts and device drivers}; Ch.~7 \emph{Locking} & -- & Sv39 tables, COW, the UART and timer, \code{amoswap} locks. \\
\must & \SPEC & Virtio: \S2.1 \emph{Device Status Field}, \S2.2 \emph{Feature Bits}, \S2.7 \emph{Split Virtqueues}; \S3.1 \emph{Device Initialization}; \S4.2 \emph{Virtio Over MMIO}; \S5.2 \emph{Block Device} & -- & All of Task 17.10. \\
\should & \SPEC & Devicetree specification \S3.4 \code{/memory}, \S3.5 \code{/reserved-memory}, \S3.6 \code{/chosen}; QEMU's \emph{`virt' Generic Virtual Platform} page for RISC-V & -- & What Tasks 17.1 and 17.6 read. \\
\should & \OSTEP & Ch.~20 \S20.3 \emph{Multi-level Page Tables} (re-read) & 215--222 & Sv32 is this design with other bit names. \\
\could & \OSTEP & Ch.~6 \emph{Mechanism: Limited Direct Execution} (re-read with Figure~\ref{fig:trap}) & 49--62 & Which part of a trap is hardware, which is yours. \\
\could & \OSTEP & Ch.~28 \S28.6--28.11 (test-and-set to fetch-and-add, with LL/SC) & 306--313 & \code{amoswap.w} is test-and-set; \code{lr.w}/\code{sc.w} are LL/SC. \\
\end{readingplan}

Keep the \emph{RISC-V ELF psABI} (\url{https://github.com/riscv-non-isa/riscv-elf-psabi-doc}) at hand for the
calling convention and the ELF flags of Task~17.11.

\begin{minixbox}[Real-system lens: how xv6-riscv does it]
xv6 runs on the same board. It has no firmware: it boots with \code{-bios none}, \code{kernel/entry.S} starts in
M-mode, and \code{start.c} delegates all traps to S-mode, opens the PMP, enables Sstc for the timer and \code{mret}s
into \code{main}; your kernel never touches an M-mode CSR. It is RV64 with Sv39, so \code{walk()} in \code{vm.c}
descends three levels to your two. Its kernel has its own page table, so \code{trampoline.S} switches \code{satp} on
every trap; Turk-OS maps the kernel everywhere and needs no trampoline. The rest maps one to one: \code{kernelvec.S}
and \code{trap.c} are your \code{entry.S} and \code{trap.c}, \code{plic.c} is a dozen lines, \code{uart.c} drives a
16550, \code{swtch.S} saves \code{ra}, \code{sp} and \code{s0}--\code{s11} (in the process structure),
\code{spinlock.c} compiles to \code{amoswap.w.aq}, and \code{virtio_disk.c} builds Task~17.10's three-descriptor request.
\end{minixbox}

\section{Build it}

\task{Meet the board}
Look before you program. Run OpenSBI alone and read its banner: \code{Domain0 Next Address} is where it will jump (0
without \code{-kernel}, the ELF entry point with one), \code{Next Arg1} the FDT it will pass, \code{Next Mode} S-mode.
Then list the memory map and dump the devicetree next to the blobs of Task~16.10.

\begin{lst}{sh}{terminal}
qemu-system-riscv32 -machine virt -nographic -bios default       # Ctrl-A X quits
qemu-system-riscv32 -machine virt -display none -monitor stdio -S # then: info mtree -f
qemu-system-riscv32 -machine virt,dumpdtb=tests/data/virt-riscv32.dtb
dtc -I dtb -O dts tests/data/virt-riscv32.dtb | less
\end{lst}

Find each row of Table~\ref{tab:virt} in the tree, and note three surprises: \texttt{\#address-cells} and
\texttt{\#size-cells} are 2 even on this 32-bit machine, so each \code{reg} number is two cells; a device's PLIC source
is its \code{interrupts} property; and there is no \code{/reserved-memory}, because OpenSBI adds it to the copy it
passes on (Task~17.6). Record it all in \code{docs/porting.md}.
\verify{Your notes give, without looking again: OpenSBI's next address for an ELF kernel, the FDT address with 128\,MiB
of RAM (\code{0x87e00000}, 2\,MiB below the end of RAM), the UART's interrupt (10) and the timebase (10\,000\,000 Hz).
The new blob passes your FDT host tests.}

\task{The first instructions}
OpenSBI enters in S-mode with \code{a0} = hart ID, \code{a1} = the FDT's physical address and \code{satp} = 0, and
promises nothing more. With OpenSBI 1.7 on QEMU 10.2, \code{sp} still points into OpenSBI's memory, which S-mode may
not touch; \code{stvec} still holds the entry point, so an early trap silently restarts the kernel; and
\code{sstatus}.FS is on, so floating-point registers would leak between processes. \code{boot.S} (GNU \code{as} through
the C preprocessor, \code{//} comments) fixes all three, zeroes \code{.bss} and calls \code{kmain}.

\begin{lst}{rvasm}{kernel/arch/riscv32/boot.S}
    .section .text.boot, "ax"
    .globl _start
_start:                         // S-mode, a0 = hart ID, a1 = FDT (physical), satp = 0
    csrw  sie, zero             // no interrupt sources until arch_init
    li    t0, 0x6000            // sstatus.FS (bits 14:13) = 0: FPU off,
    csrc  sstatus, t0           //   any floating-point instruction traps
    la    t0, early_trap        // stvec still holds OpenSBI's jump address
    csrw  stvec, t0
    la    sp, boot_stack_top    // 16 KiB in .bss, declared further down
    la    t0, _bss_start        // zero .bss (a0 and a1 untouched)
    la    t1, _bss_end
1:  bgeu  t0, t1, store_boot_args
    sw    zero, 0(t0)
    addi  t0, t0, 4
    j     1b
store_boot_args:                // boot_hartid, boot_fdt: two words in .bss,
    la    t0, boot_hartid       //   read later by arch_early_init
    sw    a0, 0(t0)
    la    t0, boot_fdt
    sw    a1, 0(t0)
    call  kmain                 // never returns
    .balign 4                   // stvec needs 4-byte alignment
early_trap:                     // stuck here? read scause and sepc in the monitor
    j     early_trap
\end{lst}

Link at \code{0x80400000}, 4\,MiB into RAM and clear of OpenSBI (QEMU puts a non-ELF kernel there too). QEMU loads
an ELF \code{-kernel} where its program headers say and passes the entry point to OpenSBI.

\begin{lst}{plain}{kernel/arch/riscv32/kernel.ld (until Task 17.7)}
OUTPUT_ARCH(riscv)
ENTRY(_start)
SECTIONS {
    . = 0x80400000;  _kernel_start = .;     /* OpenSBI jumps to our entry point */
    .text   : { *(.text.boot) *(.text .text.*) }
    .rodata : ALIGN(4K) { *(.rodata .rodata.* .srodata .srodata.*) }
    .data   : ALIGN(4K) { *(.data .data.* .sdata .sdata.*) }
    .bss    : ALIGN(4K) { _bss_start = .; *(.sbss* .bss .bss.* COMMON) . = ALIGN(16); _bss_end = .; }
    _kernel_end = .;
}
\end{lst}

Until the UART works, print through the SBI. \code{sbi.c} wraps \code{ecall} once. On top of it,
\code{sbi_putchar(c)} is DBCN function 2 (\code{console_write_byte}), or the legacy EID \code{0x01} if Base function 3
(probe) says DBCN is missing, and \code{sbi_set_timer(t)} is TIME function 0 with the 64-bit deadline split into
\code{a0} (low half) and \code{a1} (high half), as the RV32 calling convention passes a 64-bit argument. Register the
putchar with \code{console_register()} as the boot console. \code{arch_early_init} lives here too.

\begin{lst}{c}{kernel/arch/riscv32/sbi.c}
struct sbiret { long error; long value; };      /* error in a0, value in a1 */

static struct sbiret sbi_call(long eid, long fid, long arg0, long arg1, long arg2)
{
    register long a0 __asm__("a0") = arg0;
    register long a1 __asm__("a1") = arg1;
    register long a2 __asm__("a2") = arg2;
    register long a6 __asm__("a6") = fid;
    register long a7 __asm__("a7") = eid;
    __asm__ volatile ("ecall" : "+r"(a0), "+r"(a1) : "r"(a2), "r"(a6), "r"(a7) : "memory");
    return (struct sbiret){ a0, a1 };
}
\end{lst}

\begin{lst}{make}{kernel/arch/riscv32/arch.mk}
CROSS       := riscv32-elf-
ARCH_CFLAGS := -march=rv32imac_zicsr_zifencei -mabi=ilp32 -mcmodel=medany
ARCH_SRCS   := $(wildcard kernel/arch/riscv32/*.c kernel/arch/riscv32/*.S)
LDSCRIPT    := kernel/arch/riscv32/kernel.ld
QEMU        := qemu-system-riscv32
QEMU_FLAGS  := -machine virt -bios default -m 128M -nographic -kernel $(BUILD)/kernel.elf
QEMU_TEST_FLAGS  := -append "test" -no-reboot   # the sifive_test exit device is always on virt (Task 17.12)
TEST_PASS_STATUS := 0
\end{lst}
\verify{\code{make ARCH=riscv32 run} prints \code{Merhaba from S-mode: hart 0, FDT at 0x87e00000} below the OpenSBI
banner, which now reports \code{Domain0 Next Address : 0x80400000}; \code{riscv32-elf-readelf -l} shows one load
segment at that address.}

\task{The 16550 again, memory-mapped}
The console is the 16550 of Task~2.6, same registers, one byte apart, but memory-mapped; Task~16.8 already split the
driver into a core and accessors. Take the address from the FDT: follow \code{/chosen/stdout-path}, check for
\code{ns16550a} in \code{compatible}, read \code{reg} (two cells each, big-endian, or every address comes out as 0) and
\code{interrupts}. Without \code{reg-shift} the stride is one byte. Initialise it as on x86, register it as the console,
and keep the SBI console for panics.
\verify{The boot log looks the same as with the SBI console, and a GDB breakpoint on the MMIO write accessor is hit for
every character, writing to \code{0x10000000}.}

\task{Traps}
The RV32 \code{struct trapframe} holds the 31 registers in order and the four trap CSRs, padded to 144 bytes for
16-byte stack alignment.

\begin{lst}{c}{kernel/arch/riscv32/include/arch/trapframe.h}
#define SSTATUS_SIE (1u << 1)
#define SSTATUS_SPIE (1u << 5)
#define SSTATUS_SPP (1u << 8)
#define SSTATUS_SUM (1u << 18)

struct trapframe {                  /* x1..x31, then the CSRs: 36 words = 144 bytes */
    unsigned long ra, sp, gp, tp, t0, t1, t2, s0, s1, a0, a1, a2, a3, a4, a5, a6, a7;
    unsigned long s2, s3, s4, s5, s6, s7, s8, s9, s10, s11, t3, t4, t5, t6;
    unsigned long sepc, sstatus, scause, stval, pad;
};

static inline unsigned long tf_sysno(const struct trapframe *tf)      { return tf->a7; }
static inline unsigned long tf_arg(const struct trapframe *tf, int n) { return (&tf->a0)[n]; }
static inline void tf_set_ret(struct trapframe *tf, unsigned long v)  { tf->a0 = v; }
static inline uintptr_t tf_pc(const struct trapframe *tf)            { return tf->sepc; }
static inline bool tf_from_user(const struct trapframe *tf)          { return !(tf->sstatus & SSTATUS_SPP); }
/* tf_set_pc and tf_sp in the same style */
\end{lst}

\code{entry.S} needs a free register before it can save any, and \code{sp} may be a user's. The convention:
\code{sscratch} holds the top of the task's kernel stack while in U-mode, and 0 in S-mode. Swap \code{sp} and
\code{sscratch} first; a non-zero result means the trap came from U-mode, zero means swap back and stay on the kernel
stack. \code{trap_return} refills \code{sscratch} only when going back to U-mode. (\code{.irp} repeats a line for a
list.)

\begin{lst}{rvasm}{kernel/arch/riscv32/entry.S}
#define TF_SIZE 144
#define TF_SP 4
#define TF_SEPC 124
#define TF_SSTATUS 128
    .section .text
    .globl trap_entry, trap_return
    .balign 4
trap_entry:                         // stvec points here (direct mode)
    csrrw sp, sscratch, sp          // from U-mode: sp = kernel stack, sscratch = user sp
    bnez  sp, 1f
    csrrw sp, sscratch, sp          // from S-mode (sscratch was 0): undo the swap
1:  addi  sp, sp, -TF_SIZE
    .irp n, 1,3,4,5,6,7,8,9,10,11,12,13,14,15,16,17,18,19,20,21,22,23,24,25,26,27,28,29,30,31
    sw    x\n, (\n - 1) * 4(sp)     // every register but sp (x2)
    .endr
    csrrw t0, sscratch, zero        // t0 = user sp, or 0 from S-mode; sscratch = 0 now
    bnez  t0, 2f
    addi  t0, sp, TF_SIZE           // from S-mode: sp before the trap
2:  sw    t0, TF_SP(sp)
    csrr  t1, sepc
    csrr  t2, sstatus
    csrr  t3, scause
    csrr  t4, stval
    sw    t1, TF_SEPC(sp)           // sepc, sstatus, scause, stval: offsets 124-136
    sw    t2, TF_SSTATUS(sp)
    sw    t3, TF_SSTATUS + 4(sp)
    sw    t4, TF_SSTATUS + 8(sp)
    mv    a0, sp
    call  trap_handler              // void trap_handler(struct trapframe *tf)
trap_return:                        // sp = the frame to resume
    lw    t1, TF_SEPC(sp)
    csrw  sepc, t1
    lw    t1, TF_SSTATUS(sp)
    csrw  sstatus, t1               // saved SIE is 0: no interrupts from here on
    andi  t1, t1, 0x100             // SPP set: back to S-mode, sscratch stays 0
    bnez  t1, 3f
    addi  t1, sp, TF_SIZE           // to U-mode: the next trap from there
    csrw  sscratch, t1              //   starts on this kernel stack, empty again
3:  .irp n, 1,3,4,5,6,7,8,9,10,11,12,13,14,15,16,17,18,19,20,21,22,23,24,25,26,27,28,29,30,31
    lw    x\n, (\n - 1) * 4(sp)
    .endr
    lw    sp, TF_SP(sp)             // last: the user sp, or the interrupted kernel sp
    sret                            // pc = sepc, mode = SPP, SIE = SPIE
\end{lst}

\code{arch_init} sets \code{stvec} to \code{trap_entry} and \code{sscratch} to 0. In C, user exceptions become
signals; anything unexpected in S-mode gets a crash report with the cause name of Table~\ref{tab:causes}. A breakpoint
may be \code{c.ebreak} (2 bytes) or \code{ebreak} (4), and the assembler picks, so read the low bits.

\begin{lst}{c}{kernel/arch/riscv32/trap.c}
void trap_handler(struct trapframe *tf)
{
    unsigned long cause = tf->scause;
    if (cause & (1ul << 31)) {                      /* interrupt 5: timer, 9: PLIC    */
        if ((cause & 0xff) == 5) timer_tick(); else plic_dispatch();
        return;
    }
    if (tf_from_user(tf))
        csr_set(sstatus, SSTATUS_SUM);              /* Task 17.7: may touch user memory */
    if (cause == 8) {                               /* ecall from U-mode (Task 17.9)  */
        tf->sepc += 4;                              /* resume after the ecall         */
        csr_set(sstatus, SSTATUS_SIE);              /* system calls can be preempted  */
        syscall_dispatch(tf);
    } else if (cause == 3 && !tf_from_user(tf)) {   /* a kernel breakpoint            */
        uint16_t low = *(volatile uint16_t *)tf->sepc;
        tf->sepc += ((low & 3) == 3) ? 4 : 2;       /* 32-bit encodings end in 11     */
    } else if (cause == 12 || cause == 13 || cause == 15) {
        page_fault(tf);                             /* generic, via arch_fault_info   */
    } else if (tf_from_user(tf)) {
        kill_current(cause == 2 ? SIGILL : cause == 3 ? SIGTRAP : SIGSEGV);
    } else {
        crash_report(tf, cause_name(cause));       /* registers, sepc, stval; panics */
    }
}
\end{lst}
\verify{\code{__asm__ volatile ("ebreak")} in \code{kmain} prints ``breakpoint at \dots'' and the kernel carries on;
inside \code{.option norvc} the same test advances \code{sepc} by 4 instead of 2. Reading \code{*(volatile int *)0}
gives a crash report with ``load access fault'' and \code{stval=00000000}: with paging off, address 0 is no memory.}

\task{Time and interrupts}
Read \code{timebase-frequency} from \code{/cpus}; \code{arch_timer_init(hz)} computes the interval (10\,000\,000 /
100 = 100\,000 counts), calls \code{sbi_set_timer(now + interval)} and sets \code{sie}.STIE (bit 5). On RV32 the 64-bit
\code{time} counter comes in two halves: read \code{rdtimeh}, \code{rdtime}, \code{rdtimeh} again, and retry if the high
half changed. \code{arch_cycles} returns that value (timebase counts, not CPU cycles). \code{timer_tick} adds the
interval to the \emph{previous} deadline, so the tick does not drift, calls \code{sbi_set_timer} again (without it the
timer fires exactly once) and then the generic tick code of Phases 4 and 8.

On QEMU \code{virt} the PLIC context for hart 0 in S-mode is context 1, as the PLIC node's
\code{interrupts-extended} property says. \code{plic_init(base)} maps three of the PLIC's pages with \code{devmap()}
(Task~17.7; until then it returns the physical address): the priorities (source $i$ at base $+4i$), the enable bits of
context 1 (base \code{+ 0x2080}), and its threshold and claim registers (base \code{+ 0x201000} and \code{+ 0x201004}).
It sets the threshold to 0 and \code{sie}.SEIE (bit 9).

\code{arch_irq_unmask(irq)} sets the source's priority to 1 (0 means never) and its enable bit; the dispatcher
claims, calls, completes:

\begin{lst}{c}{kernel/arch/riscv32/plic.c}
void plic_dispatch(void)                            /* scause = interrupt 9           */
{
    uint32_t irq;
    while ((irq = ctx[1]) != 0) {                   /* claim (base + 0x201004)        */
        if (irq_handlers[irq])
            irq_handlers[irq](irq);
        ctx[1] = irq;                               /* complete, or it never returns  */
    }
}
\end{lst}

Input comes from the UART now. Set the 16550's interrupt enable register to 1 (data received), register a handler
for source 10 and unmask it; the handler in \code{uart16550.c} passes bytes to \code{tty_input_char()} while bit 0 of
the line status register is set. With \code{sstatus}.SIE set in \code{arch_init}, Task~4.6's monitor runs over serial.
\verify{\code{uptime} in the monitor advances by one second per second. Typed characters arrive by interrupt (a
breakpoint in the UART handler is hit), and nothing polls the UART.}

\task{Memory from the devicetree}
Fill \code{struct bootinfo} from the FDT that OpenSBI passes in \code{a1}, not the dumped one: OpenSBI adds
\code{/reserved-memory} nodes (\code{no-map}) for its 320\,KiB at \code{0x80000000}. Phase 5's allocator is unchanged.

Usable RAM is the \code{reg} of each \code{memory@...} node. Reserved are the children of \code{/reserved-memory},
the FDT header's reservation block, the FDT itself (\code{totalsize} bytes) and the kernel image
(\code{_kernel_start}--\code{_kernel_end}). \code{/chosen} gives the command line (\code{bootargs}, from
\code{-append}), the initrd (\code{linux,initrd-start} and \code{-end}, from \code{-initrd}), the console
(\code{stdout-path}) and \code{rng-seed}, a seed for \code{arch_entropy}.

To capture what the kernel receives, use the QEMU monitor: \code{pmemsave 0x87e00000 0x10000 "passed.dtb"}, and add
the blob to the FDT tests.
\verify{\code{mem} reports 128\,MiB (256\,MiB with \code{-m 256M}) and lists \code{0x80000000}--\code{0x8004ffff}, the
kernel and the FDT as reserved. A test that allocates every free frame finds none in a reserved range, and the Phase 5
allocator tests pass.}

\task{Sv32 and the higher half}
Now move. In \code{<arch/memlayout.h>}, set \code{KERNEL_BASE} to \code{0xC0000000} and keep \code{PHYS_BASE} at
\code{0x80000000}, so \code{P2V(pa)} is \code{pa - PHYS_BASE + KERNEL_BASE}: RAM appears at \code{0xC0000000}, the
kernel at \code{0xC0400000}, and the direct map may cover 512\,MiB as on x86. In \code{kernel.ld}, start at
\code{0xC0400000}, give every output section a load address \code{AT(ADDR(...) - 0x40000000)}, and make the entry point
physical (\code{ENTRY(_start_phys)} with \code{_start_phys = _start - 0x40000000}), because OpenSBI jumps there with
translation off.

\code{boot.S} still runs at physical addresses, which works because \code{la} is pc-relative. A static table maps
RAM with 4\,MiB megapages twice: root entries \code{0x200}--\code{0x27F} map \code{0x80000000}+ to itself (only until
the jump) and \code{0x300}--\code{0x37F} map \code{0xC0000000}+ to \code{0x80000000}+, 512\,MiB each. GNU \code{as}
builds such a table at assembly time with \code{.rept 128}, \code{.word ((pa >> 12) << 10) | 0xCF} (V, R, W, X, A, D) and
\code{.set pa, pa + 0x400000}; \code{.zero} fills the gaps. Then:

\begin{lst}{rvasm}{kernel/arch/riscv32/boot.S (paging part)}
    .section .text.boot, "ax"
enable_paging:                          // (first: boot_pgdir[0x3FF] -> dev_pt, see below)
    la    t1, boot_pgdir                // physical, since la is pc-relative
    srli  t1, t1, 12
    li    t2, 1 << 31                   // satp.MODE = Sv32
    or    t1, t1, t2
    csrw  satp, t1                      // translation on; the identity entries
    sfence.vma                          //   keep the next fetch working
    lui   t0, %hi(higher_half)          // the absolute link address 0xC04...
    addi  t0, t0, %lo(higher_half)
    jr    t0
higher_half:                            // next: stvec = trap_entry (it held a physical
    la    sp, boot_stack_top            //   address), zero the 128 identity entries,
    j     store_boot_args               //   sfence.vma; then Task 17.2's code and kmain
\end{lst}

Devices live in the last 4\,MiB: root entry \code{0x3FF} points to a page table \code{dev_pt} in \code{.bss}, and
\code{devmap(pa)} in \code{sv32.c} fills its next entry with a read-write leaf (A and D set) and returns the address. On
RISC-V, whether an address is RAM or a device is fixed by the platform, so Sv32 has no cache bits like x86's PCD.

The rest of \code{sv32.c} is Phase 6's VMM in RISC-V notation, with \code{PTE_PPN(pa)} = \code{((pa >> 12) << 10)}
and \code{PTE_COW} = RSW bit 8:

\begin{lst}{c}{kernel/arch/riscv32/sv32.c}
static pte_t vm_to_pte(unsigned flags)
{
    pte_t p = PTE_V | PTE_A;                  /* A and D set by us: no fault to set them */
    if (flags & VM_READ)  p |= PTE_R;
    if (flags & VM_WRITE) p |= PTE_R | PTE_W | PTE_D;     /* W without R is reserved */
    if (flags & VM_EXEC)  p |= PTE_X;
    if (flags & VM_USER)  p |= PTE_U;
    if (flags & VM_COW)   p = (p & ~PTE_W) | PTE_COW;     /* read-only until written */
    return p;
}

static pte_t *walk(pte_t *root, uintptr_t va, bool create)
{
    pte_t *pde = &root[va >> 22];                         /* VPN[1] */
    if (!(*pde & PTE_V)) {
        paddr_t pt = create ? pmm_alloc_frame() : 0;
        if (!pt)
            return NULL;
        memset(P2V(pt), 0, PAGE_SIZE);
        *pde = PTE_PPN(pt) | PTE_V;                       /* R = W = X = 0: a pointer */
    }
    ASSERT(!(*pde & (PTE_R | PTE_W | PTE_X)));            /* not inside a megapage    */
    pte_t *table = P2V(PTE_ADDR(*pde));
    return &table[(va >> 12) & 0x3FF];                    /* VPN[0] */
}
\end{lst}

\code{arch_mmu_switch} writes \code{(1 << 31) | (V2P(root) >> 12)} to \code{satp}, then \code{sfence.vma zero,
zero}; \code{arch_tlb_flush(va)} is \code{sfence.vma va, zero}; map, unmap, lookup and protect use \code{walk}.
\code{arch_mmu_init_space} copies root entries 768--1023 from \code{boot_pgdir}, now the kernel's address space, so
pre-allocate the kernel tables as in Task~6.6. \code{arch_fault_info} takes \code{addr} from \code{stval}, \code{write}
from cause 15, \code{exec} from 12, \code{user} from SPP, and \code{present} from a lookup. Task~11.4's copy-on-write then
works unchanged.

With \code{sstatus}.SUM clear, any S-mode access to a U page faults, \code{copy_from_user} included. Turk-OS sets SUM
for traps from U-mode (in \code{trap.c}); the saved \code{sstatus} has it clear, so \code{sret} clears it again.
\code{arch_signal_frame} and \code{arch_sigreturn} set it themselves.
\verify{\code{info registers} shows the pc at \code{0xc04...}; the monitor's \code{info mem} lists \code{c0000000} with
\code{rwx--ad}, the device pages at \code{ffc00000} and nothing below \code{0xC0000000}. Reading address 0 now gives a
``load page fault'', and the VMM tests of Task~6.7 pass.}

\task{Tasks and locks}
\code{arch_switch} is Phase 8's \code{switch_context}: save \code{ra} and \code{s0}--\code{s11} on the old stack,
swap stack pointers, restore.

\begin{lst}{rvasm}{kernel/arch/riscv32/switch.S}
// void arch_switch(struct context **save, struct context *load)
    .section .text
    .globl arch_switch, task_trampoline
arch_switch:
    addi  sp, sp, -64               // 13 words, rounded up for 16-byte alignment
    .set  off, 0
    .irp  r, ra, s0, s1, s2, s3, s4, s5, s6, s7, s8, s9, s10, s11
    sw    \r, off(sp)
    .set  off, off + 4
    .endr
    sw    sp, 0(a0)                 // *save = sp
    mv    sp, a1                    // the new task's kernel stack
    .set  off, 0
    .irp  r, ra, s0, s1, s2, s3, s4, s5, s6, s7, s8, s9, s10, s11
    lw    \r, off(sp)
    .set  off, off + 4
    .endr
    addi  sp, sp, 64
    ret                             // to where the new task called arch_switch
task_trampoline:                    // ra of a new task (arch_task_init)
    csrsi sstatus, 2                // SIE = 1: Phase 8's sti
    mv    a0, s1                    // entry(arg), with s0 = entry, s1 = arg
    jalr  s0
    call  task_exit                 // never returns
\end{lst}

\code{arch_task_init} fakes a 64-byte context (\code{ra} = \code{task_trampoline}, \code{s0} = entry, \code{s1} =
argument) as in Task~8.4. \code{irq_save} is \code{csrrci} on SIE, which clears it and returns the old \code{sstatus};
\code{cpu_idle} is \code{wfi}, which wakes for an enabled pending interrupt even while SIE is 0.

\begin{lst}{c}{kernel/arch/riscv32/include/arch/atomic.h}
#ifdef __riscv_atomic                   /* defined by the compiler when 'a' is in -march */
#define ARCH_HAS_ATOMICS 1
static inline uint32_t atomic_xchg(volatile uint32_t *p, uint32_t v)
{
    uint32_t old;
    __asm__ volatile ("amoswap.w.aq %0, %2, %1" : "=r"(old), "+A"(*p) : "r"(v) : "memory");
    return old;
}
#else
#define ARCH_HAS_ATOMICS 0              /* one hart, no A extension (Task 17.13) */
#endif
\end{lst}

\code{amoswap.w.aq} is Task~9.2's test-and-set with acquire ordering; release with \code{__atomic_store_n} and
\code{__ATOMIC_RELEASE}. (\code{lr.w}/\code{sc.w} serve more complex updates.) With \code{ARCH_HAS_ATOMICS} 0 a spinlock
is only \code{irq_save}, which on one hart is mutual exclusion; there, avoid \code{__atomic} builtins, which become calls
such as \code{__atomic_exchange_4} that no kernel has.
\verify{Phase 8's demo threads preempt each other, the Phase 9 tests (bounded buffer, dining philosophers) pass, and
1000 create/exit cycles leave the frame count unchanged.}

\task{Into user mode and back}
\code{arch_user_entry} prepares a frame for \code{sret}; \code{arch_return_to_user} (\code{user.S}) clears SIE, sets
\code{sp} to the frame at the top of the kernel stack and jumps to \code{trap_return}.

\begin{lst}{c}{kernel/arch/riscv32/trap.c (user entry)}
void arch_user_entry(struct trapframe *tf, uintptr_t entry, uintptr_t sp)
{
    memset(tf, 0, sizeof *tf);              /* no kernel values in user registers */
    tf->sepc = entry;                       /* sret jumps here                    */
    tf->sp = sp;                            /* argc, argv, envp (Task 11.6)       */
    unsigned long s = csr_read(sstatus);
    s &= ~(SSTATUS_SPP | SSTATUS_SIE | SSTATUS_SUM | SSTATUS_FS);  /* SPP = 0: U */
    tf->sstatus = s | SSTATUS_SPIE;         /* interrupts on after sret           */
}
\end{lst}

A system call is cause 8; \code{sepc += 4} comes \emph{before} \code{syscall_dispatch}, because \code{execve} and
\code{sigreturn} replace \code{sepc}, and \code{fork} copies the frame. Arguments come from \code{a0}--\code{a5}, the
number from \code{a7}: still the Linux i386 number (\code{write} is 4, not Linux RISC-V's 64). Task~10.9's evil programs
give two new answers. \code{divzero} is \emph{not} killed (its quotient is $-1$; print it), and \code{cli} becomes
\code{csrci sstatus, 2}, killed as an illegal instruction (2). \code{nullptr} and \code{kpeek} die of a load page fault
(13); \code{badptr}, \code{hugelen} and \code{badsys} get \code{-EFAULT}, \code{-EFAULT} and \code{-ENOSYS} as on x86.
\verify{An RV32 \code{hello} prints ``Merhaba from U-mode, pid 3'' and exits, a looping program is preempted, and every
evil program behaves as in the table in one boot.}

\task{A disk on virtio}
The disk is a virtio block device (Concepts aside: a virtio device shares rings of buffer descriptions with the
driver instead of imitating old hardware). QEMU's virtio-mmio defaults to the legacy interface (version 1); ask for
version 2:

\begin{lst}{sh}{kernel/arch/riscv32/arch.mk (QEMU_FLAGS, continued)}
-global virtio-mmio.force-legacy=false \
-drive file=$(BUILD)/turkfs.img,if=none,format=raw,id=hd0 \
-device virtio-blk-device,drive=hd0,bus=virtio-mmio-bus.0       # slot 0x10001000, IRQ 1
\end{lst}

Keep \code{virtio_mmio.c} and \code{virtio_blk.c} free of RISC-V, since Phase 18 reuses them. The arch code offers
each \code{virtio,mmio} node, mapped with \code{devmap}; empty slots have device ID 0. The \code{VIRTIO_MMIO_*} offsets
come from the specification's \emph{Virtio Over MMIO}: magic \code{0x000}, version \code{0x004}, device ID
\code{0x008}, features \code{0x010}--\code{0x024}, queue registers \code{0x030}--\code{0x044}, notify \code{0x050},
interrupt status and ACK \code{0x060}/\code{0x064}, status \code{0x070}, queue addresses \code{0x080}/\code{0x090}/\code{0x0a0},
configuration from \code{0x100}.

\begin{lst}{c}{kernel/drivers/virtio_mmio.c (probe and handshake)}
#define R(off) regs[(off) / 4]

int virtio_mmio_setup(volatile uint32_t *regs, uint32_t device_id, struct virtq *q)
{
    if (R(VIRTIO_MMIO_MAGIC) != 0x74726976 || R(VIRTIO_MMIO_DEVICE_ID) != device_id)
        return -ENODEV;                           /* not "virt", or another device */
    if (R(VIRTIO_MMIO_VERSION) != 2)
        return -ENODEV;                           /* legacy: see force-legacy above */
    R(VIRTIO_MMIO_STATUS) = 0;                    /* reset */
    R(VIRTIO_MMIO_STATUS) = ST_ACKNOWLEDGE;
    R(VIRTIO_MMIO_STATUS) |= ST_DRIVER;
    R(VIRTIO_MMIO_DEV_FEAT_SEL) = 1;              /* feature bits 32-63 */
    if (!(R(VIRTIO_MMIO_DEV_FEAT) & 1))           /* bit 32: VIRTIO_F_VERSION_1 */
        return -ENODEV;
    R(VIRTIO_MMIO_DRV_FEAT_SEL) = 1;  R(VIRTIO_MMIO_DRV_FEAT) = 1;   /* accept only it */
    R(VIRTIO_MMIO_DRV_FEAT_SEL) = 0;  R(VIRTIO_MMIO_DRV_FEAT) = 0;
    R(VIRTIO_MMIO_STATUS) |= ST_FEATURES_OK;
    if (!(R(VIRTIO_MMIO_STATUS) & ST_FEATURES_OK))
        return -ENODEV;                           /* the device refused */
    R(VIRTIO_MMIO_QUEUE_SEL) = 0;                 /* queue 0; check QueueNumMax first */
    R(VIRTIO_MMIO_QUEUE_NUM) = QSIZE;             /* e.g. 8 */
    set_addr(regs, VIRTIO_MMIO_QUEUE_DESC_LOW, q->desc);     /* physical: low, high */
    set_addr(regs, VIRTIO_MMIO_QUEUE_DRIVER_LOW, q->avail);
    set_addr(regs, VIRTIO_MMIO_QUEUE_DEVICE_LOW, q->used);
    R(VIRTIO_MMIO_QUEUE_READY) = 1;
    R(VIRTIO_MMIO_STATUS) |= ST_DRIVER_OK;        /* live from here on */
    return 0;
}
\end{lst}

The split virtqueue has three parts, a page each from \code{kalloc_page}: the descriptor table (buffer address,
length, flags NEXT and WRITE, next index), the available ring (requests from the driver) and the used ring (completions
from the device). A request chains three descriptors: the header \texttt{struct virtio\_blk\_req \{ uint32\_t type,
reserved; uint64\_t sector; \}} (type 0 read, 1 write; NEXT), the buffer (NEXT, plus WRITE for a read) and a status
byte (WRITE; 0 = OK). Post the head in the available ring, fence (\code{__sync_synchronize()}), bump its \code{idx},
fence, write 0 to \code{QueueNotify}, and sleep. The handler (PLIC source 1) acknowledges \code{InterruptStatus}, walks
the used ring and wakes the waiter. Sectors are 512 bytes. Register \code{vda} through Task~12.5's
\code{struct blockdev}, and the buffer cache and TurkFS sit on top unchanged.

\begin{warnbox}[Physical addresses for every buffer]
\code{V2P} is only right for direct-map addresses. A 1\,KiB buffer from the Phase 7 heap can straddle two pages
that are not adjacent in RAM, and the device writes half of it over someone else's memory. Take buffer data from
\code{kalloc_page} (four per page), or split such a buffer into two descriptors.
\end{warnbox}
\verify{The kernel prints \code{vda: virtio-blk, 131072 sectors (64 MiB)} (from the \code{capacity} field at
\code{0x100}); the Phase 12 tests pass on \code{vda}; a pattern written to sector 100 shows on the host with
\code{hexdump -C -s $((100*512)) -n 64 build/riscv32/disk.img}.}

\task{Userland for RV32}
Turk-libc and the utilities do not change. The system call stub sets \code{a7} and \code{a0}--\code{a5}, executes
\code{ecall} and returns \code{a0}, with register variables like \code{sbi_call}; \code{user.ld} still starts at
\code{0x00400000}; and \code{crt0.S} reads the System V stack of Task~11.6:

\begin{lst}{rvasm}{user/arch/riscv32/crt0.S}
    .section .text
    .globl _start
_start:                             // sp -> argc, argv[], NULL, envp[], NULL
    lw    a0, 0(sp)                 // argc
    addi  a1, sp, 4                 // argv
    slli  t0, a0, 2
    add   a2, a1, t0
    addi  a2, a2, 4                 // envp = argv + argc + 1
    andi  sp, sp, -16               // calls need a 16-byte aligned sp
    call  main                      // main(argc, argv, envp)
    call  exit                      // exit(main's result), already in a0
\end{lst}

Build with the kernel's \code{-march}/\code{-mabi}, \code{-nostdlib -static} and \code{-lgcc} (RV32 has no 64-bit
divide; GCC calls \code{__udivdi3}). The ELF loader expects \code{EM_RISCV} (243), \code{ELFCLASS32}, little-endian,
and \code{e_flags & 0x0006} = 0 (soft float: the kernel saves no FPU registers); bit \code{0x0001} (compressed) is fine.
Rebuild the root archive (Task~13.4) and the TurkFS image, and pass the archive with \code{-initrd}.
\verify{\code{make ARCH=riscv32 run} boots to \code{turk@os:/$}; \code{ls /bin}, \code{echo merhaba | cat}, a
redirection and Ctrl+C on \code{cat} behave as on x86; an i386 ELF fails \code{execve} with \code{-ENOEXEC}.}

\task{make test on RISC-V}
The test device replaces \code{isa-debug-exit}: a 32-bit write of \code{0x5555} ends QEMU with status 0, and
\code{(code << 16) | 0x3333} with status \code{code} (\code{0x7777} resets). So a pass is 0 here, not 1: set \code{TEST_PASS_STATUS := 0} in
\code{arch.mk}, and fail with a non-zero code, since code 0 also exits with 0. \code{arch_halt} is SRST function 0 with reset type 0 (shutdown). \code{-append "test"} reaches
\code{/chosen/bootargs}, so the test target stays generic. Fuzz for an hour (new code sits under the system call
boundary), then run Task~15.5's benchmarks and put them beside x86 in \code{docs/porting.md}. User-space \code{rdtime}
needs \code{scounteren}.TM (bit 1); one count is 100\,ns. Both machines are emulated on one host, so compare ratios, not
absolute times.
\verify{\code{make ARCH=riscv32 test} prints \code{ALL TESTS PASSED}, a deliberately broken test makes it fail with
status 1, \code{make ARCH=i386 test} still passes, and an hour of fuzzing ends without a panic.}

\task{Optional: Turk-OS on Timur-RV32IMC}
This depends on your hardware project: start when Timur passes its own tests and runs a bare-metal program on its
UART. Check what the core offers, fill the gaps in firmware, and build a kernel that needs nothing more.

\begin{center}
\small\renewcommand{\arraystretch}{1.15}
\begin{tabularx}{\linewidth}{@{}>{\raggedright\arraybackslash}p{0.25\linewidth} >{\raggedright\arraybackslash}X >{\raggedright\arraybackslash}p{0.33\linewidth}@{}}
\toprule
\tkth{Requirement} & \tkth{Why Turk-OS needs it} & \tkth{Workaround if missing} \\
\midrule
Zicsr; M, S, U modes; Sv32 with \code{sfence.vma} & traps; protection; address spaces, COW, the higher half & none \\
Precise traps, \code{medeleg}/\code{mideleg} & faults restart the instruction; S-mode handles them & an M-mode handler forwarding to \code{stvec} \\
\code{mtime}, \code{mtimecmp} & the timer behind \code{sbi_set_timer} & none: the tick is essential \\
A 16550-like UART, an interrupt controller & console; input by interrupt & another driver; poll the UART on each tick \\
A extension & \code{amoswap} spinlocks & build \code{rv32imc}, \code{ARCH_HAS_ATOMICS=0} \\
Misaligned access & compiled code may assume it & \code{-mstrict-align} or M-mode emulation \\
A boot loader and a devicetree & kernel, initrd and FDT in RAM; all boot facts & UART or JTAG loading; a hand-written \code{timur.dts} for \code{dtc} \\
\bottomrule
\end{tabularx}
\end{center}

Something must run in M-mode. Either port OpenSBI, whose \emph{generic} platform is driven by the FDT, or write a
small SBI shim: an \code{mtvec} handler for Base probes, TIME (program \code{mtimecmp}; on the machine timer interrupt
set \code{mip}.STIP and mask MTIE), DBCN and SRST, with everything else delegated through \code{medeleg} and
\code{mideleg}, ending in an \code{mret} to the kernel with \code{a0} = hart ID and \code{a1} = FDT. Develop it on QEMU:
\code{-bios build/shim.elf -kernel build/riscv32/kernel.elf} runs your shim in M-mode at \code{0x80000000} instead of
OpenSBI (\code{-bios none} boots one image in M-mode, as xv6 does). Build the Timur kernel with \code{-march=rv32imc_zicsr_zifencei},
an initrd-only root and the UART console, and fill the table with Timur's real answers in \code{docs/porting.md}.
\verify{Your shim boots the unchanged QEMU kernel to the shell, and the \code{rv32imc} kernel contains no
\code{amoswap}, \code{lr.w} or \code{sc.w} (\code{riscv32-elf-objdump -d}). On Timur, the milestone is the banner, then
the prompt.}

\section{Testing and debugging}

\code{-d int,guest_errors -D build/riscv32/qemu.log} logs every trap with its cause, \code{sepc} and \code{stval}.
With \code{-s -S}, a GDB that knows RISC-V (Task~16.1) and \code{set architecture riscv:rv32} steps through
\code{boot.S}. The monitor's \code{info registers} shows every CSR and \code{info mem} the live Sv32 mappings.

\begin{pitfalls}
\code{unrecognized opcode `csrr a0,sstatus', extension `zicsr' required} & Current binutils keep the CSR instructions out of the base ISA & Add \code{_zicsr_zifencei} to \code{-march}. \\
OpenSBI's banner, then nothing & Linked at the wrong address, or a virtual entry point & \code{readelf -l}; the banner's \code{Domain0 Next Address} must be your entry. \\
The banner repeats forever & A trap before \code{stvec} was set, or a call on OpenSBI's \code{sp} & Set \code{stvec} and \code{sp} first in \code{boot.S}. \\
\code{ebreak} or \code{ecall} loops, or an instruction is skipped & \code{sepc} not advanced, or by 4 over a 2-byte \code{c.ebreak} & \code{ecall}: +4 before dispatching; \code{ebreak}: low bits \code{11} mean 4 bytes. \\
Load page fault in \code{copy_from_user} & \code{sstatus}.SUM clear & Set SUM for traps from U-mode. \\
Fault right after the \code{satp} write & The pc is not mapped in the new table, or no \code{sfence.vma} & Keep the identity megapages until the jump. \\
Every first access to a page faults, on a real core & A or D clear, and the core faults instead of setting them & Set A always, D on writable leaves. \\
PLIC interrupts never arrive & Priority 0, the wrong context, SEIE clear, or a claim never completed & Context 1; priority 1, threshold 0; always complete. \\
The virtio probe sees version 1 & The legacy interface & \code{-global virtio-mmio.force-legacy=false}. \\
Corruption near \code{0x80000000} & OpenSBI's memory given out as free & Reserve \code{/reserved-memory} and the FDT. \\
\end{pitfalls}

\section{Milestone}

\begin{milestonebox}[Merhaba, RISC-V]
\begin{checklist}
  \item \code{qemu-system-riscv32 -machine virt} boots Turk-OS on OpenSBI to the \code{tsh} prompt; input arrives by
    UART interrupt.
  \item The kernel runs at \code{0xC0400000} under Sv32, and every device comes from the devicetree.
  \item Every exception gives a readable crash report, \code{ebreak} returns, and the evil programs are contained.
  \item \code{make ARCH=riscv32 test} passes, \code{make ARCH=i386 test} still passes, and an hour of fuzzing ends
    without a panic.
  \item A TurkFS image written by Turk-OS on RISC-V passes \code{fsck.minix -f} on the host.
  \item \code{docs/porting.md} records the board, every leak you fixed, and the benchmarks beside x86.
  \item (Optional) The Timur-RV32IMC table is filled in, and your SBI shim boots the kernel in QEMU.
\end{checklist}
\end{milestonebox}

\begin{questionbox}
\begin{enumerate}
  \item What can M-mode do that S-mode cannot? Which jobs of the BIOS and GRUB does OpenSBI take over?
  \item What does the hardware do on a trap, and what must \code{entry.S} do because the hardware doesn't?
  \item Why does \code{ecall} need \code{sepc += 4} while a page fault must leave \code{sepc} alone? And \code{ebreak}?
  \item How do an x86 PTE and an Sv32 PTE tell a table from a page, and ``absent'' from ``present but forbidden''?
  \item What is \code{sstatus}.SUM for? Why does Turk-OS set A and D itself when QEMU would?
  \item Why can an RV32IMC core without the A extension still run Turk-OS on one hart, and what breaks with two?
\end{enumerate}
\end{questionbox}

\section{Going further}

Port to \textbf{RV64} with Sv39, like xv6. Start more harts with the SBI \textbf{HSM} extension and make the scheduler
SMP-safe. Use \textbf{Sstc} (\code{stimecmp}, which QEMU's hart has) to drop the SBI call from every tick. Add global
kernel mappings and ASIDs. Or run Turk-OS on a real RV64 board, whose boot chain (U-Boot SPL, OpenSBI, U-Boot) is the
one QEMU's \code{virt} imitates.

\nextphase{18}{The ARM Port}
\booklegend
\portlegend
\end{document}
